\section{ELECTRIQUE}

\subsection{Description des manipulations}

Nous utilisons un cable coaxial de type RG58 de 50 $\ohm$ et de 100 mètres dans lequel nous faisons passer grace a un generateur, une impulsion de 210 ns a une frequence de 30 kHz (voir Fig. \ref{montage_elec}).
%\begin{figure}[ht]
%	\centering
%	\includegraphics[scale=0.4]{montage_elec.png}
%        \caption{Montage experience cable coaxial}
%        \label{fig:montage_elec}
%\end{figure}

\vspace{1cm}
\noindent \begin{minipage}{0.4\textwidth}
\includegraphics[width=\textwidth]{montage_elec.png}
	\captionof{figure}{U$^2$(1/D$^2$)}
\label{fig:montage_elec}
\end{minipage}
\hspace{1cm}
\begin{minipage}{0.5\textwidth}
Nous mesurons la difference entre les signaux du generateur et ceux passant par le cable. Cela nous permet de connaitre le temps de propagation de l'onde electrique dans le cablem.Si on connait le temps et la distance, nous pouvons ainsi connaitre la vitesse de propagation grace a la relation Eq. 2.

\noindent \end{minipage}

\vspace{1cm}


En un second temps, nous utilisons un pont RLC (voir Fig. \ref{fig:rlc}) afin de determiner la capacite lineique c* et l'inductance lineique l* en y connectant une extremite de notre cable coaxial. 

\vspace{1cm}
\noindent \begin{minipage}{0.4\textwidth}
\includegraphics[width=\textwidth]{rlc.png}
        \captionof{figure}{Montage d'un pont RLC}
\label{fig:rlc}
\end{minipage}
\hspace{1cm}
\begin{minipage}{0.5\textwidth}
	Afin de determiner c*, le pont RLC est utile dans la mesure ou, en laissant l'autre extremite du cable ouvert, le cable devient une capacite. A l'inverse, pour trouver l*, nous fermons le circuit. Ainsi notre cable se comporte comme une inductance.  
\noindent \end{minipage}

\vspace{1cm}
Enfin, en faisant varier les resistances de terminaison (Voir R sur Fig. \ref{fig:montage_elec}), nous pouvons mesurer les amplitudes des impulsions transmises et reflechies afin de trouver les coefficients de reflexion et de transmission.
